\section{Lead9 contact-perimeter field attribution (V3.15)}
\label{sec:v315-perimeter}
The source-native 3D lead9/solder electrical contacts remain the limiting raw-current-gradient problem even after the separate V3.14 fine-mesh P2 solve. This study uses \emph{only the two archived P2 fields}, their unchanged original 12-region source geometry, and the existing 11 finite-contact ownership rules. No new PDE is assembled or solved. On each physical contact facet the two raw currents come from \(-\boldsymbol\sigma\nabla V\), whereas the distinct finite-contact-law current is evaluated from \((V_1-V_2)/r''_e\) using the saved material-separated quadratic traces. The original contact areal resistivity hypothesis is \(r''_e=10^{-11}\ \Omega\mathrm m^2\), not an OEM measurement.

\fig[.88]{v315_lead9_rim_mismatch}{pdf}{Source-native radial TRI3 centroid study of the solder-side absolute raw-gradient/contact-law mismatch. The annular diagnostic bin \(0.90\leq r/r_{\rm CAD}<1.05\) covers about one-fifth of each registered contact surface. The ordinate is the \emph{fraction of that contact's absolute total mismatch} associated with that radial bin, not the local current or a prediction of a newly modified contact.}

For lead9 to the BiSn land, this outer band contains \textbf{78.01\%} of the earlier P2 solder-side absolute mismatch and \textbf{79.11\%} after refinement; for lead9 to the InSn pin joint the corresponding fractions are \textbf{72.47\%} and \textbf{75.94\%}. The registered rim area fractions remain between 19.2\% and 20.4\%. Meanwhile the full-interface raw-current cut errors reduce from 8.0993\% to 7.0193\% (land) and 6.8529\% to 6.1887\% (pin), still greatly exceeding the unchanged \textbf{0.1\%} release limit.

\textbf{Frozen-field negative control.} A counterfactual that zeroes \emph{all of the existing rim's signed raw-minus-contact-law bias} but leaves the interior solved field unchanged would still yield worst signed-side errors of \textbf{2.123\%} at the land and \textbf{1.304\%} at the pin for the finer P2 solution. The counterfactual is not a PDE solution or error bound for a newly solved field, but establishes that removing the currently observed rim discrepancy \emph{alone} is not sufficient. The original physical electrical Gate B remains \textbf{FAIL/STOP}; thermal Gate C and coupled Joule Gate D remain \textbf{NOT RUN}. Future work should favor a separately registered contact-edge/interface-formulation study with original raw-gradient comparisons rather than a generic mesh-size sweep. The complete numerical decomposition and source hashes are documented in \texttt{docs/V3\_15\_LEAD9\_CONTACT\_PERIMETER\_FIELD\_DIAGNOSIS\_20261110.md}.
